\label{chap:Datasets}
This chapter describes the case study and the \ac{ML} pre-trained model and
dataset used in this study for detecting the arrival times of primary (P) and
secondary (S) waves. A summary of the benchmark datasets used to train PhaseNet 
(\cite{PhaseNet}) and EQTransformer (\cite{EQTransformer}) is provided in Table 
\ref{tab:datasets}. These datasets differ in geographical coverage, noise 
characteristics, and labeling quality, offering complementary advantages for 
benchmarking and evaluation.

\section{Case Study}
\label{sec:case_study}

The primary dataset used in this thesis consists of continuous seismic 
recordings from the \ac{SMINO}, managed by the \ac{OGS} and surrounding 
seismic network. This dataset serves as the target domain for evaluating the 
efficacy of the implemented machine-learning pipeline in a operational setting.

Key characteristics include:
\begin{itemize}
  \item \textbf{Geographical setting}: Northeastern Italy, with a focus on the 
  Friuli Venezia Giulia (FVG) and Veneto regions. This area is characterized by 
  complex tectonics resulting from the collision between the Adria microplate 
  and the Eurasian plate, leading to significant seismic hazard.
  \item \textbf{Instrumentation}: 266 stations across 18 networks
  (Fig.~\ref{fig:OGSStations}), comprising a mix of high-sensitivity broadband
  seismometers and strong-motion accelerometers. This dense network ensures
  high-fidelity recording of both microseismicity and larger regional events.
  \item \textbf{Sampling rate}: All waveforms are resampled to 100 Hz to align 
  with the input specifications of the pre-trained deep-learning models.
  \item \textbf{Observation period}: A three-month continuous period from 
  March 20, 2024, to June 20, 2024. This timeframe captures a representative 
  range of seismic activity (Fig.~\ref{fig:OGSSeismicity}), The period under 
  analysis was chosen because on 27 March 2024 a magnitude 4.6 earthquake 
  occurred in Friuli, near Socchieve (Lat: 46.3583, Lon: 12.808; Depth: 10.22 
  km; UTC: 2024-03-27 21:19:37; Local time: 22:19:37; Ml: 4.6 ± 0.3; Location: 
  5 km SSO of Socchieve, Udine). This is one of the most energetic events to 
  occur in Friuli since 1980.
  \item \textbf{Format}: Data are retrieved in continuous miniSEED (MSEED) 
  format directly from multiple data centers (OGS, INGV, IRIS, GFZ, ETH,
  ORFEUS, and GEOFON) using the ObsPy library (\cite{ObsPy}), with station 
  availability shown in Fig.~\ref{fig:OGSAvailability}. Metadata are managed 
  using StationXML, ensuring accurate station coordinates and instrument 
  response information. To compare the ML results with the OGS reference 
  catalog, OGS events were relocated using the 1D NonLinLoc (\cite{NonLinLoc}) 
  configuration adopted in this study. The catalog will be published soon 
  (Magrin, personal communication). It contains 380 events within the red 
  polygon shown in Figure~\ref{fig:OGSSeismicity}, with local magnitudes 
  ranging from -0.2 to 4.6. A total of 7,129 phase picks were associated with 
  these events, including 4,126~P picks and 3,003~S picks.
\end{itemize}

\begin{figure}[h!]
  \centering
  \includegraphics[width=\textwidth]{OGSCatalog/GaMMA0.1/img/OGSStations.png}
  \caption{Geographical distribution of seismic stations used in this study 
  (see Data and resources section for details).}
  \label{fig:OGSStations}
\end{figure}

\begin{figure}[h!]
  \centering
  \includegraphics[width=\textwidth]{OGSCatalog/GaMMA0.1/img/.all_EventsMap.pdf}
  \caption{Seismicity recorded during the observation period.}
  \label{fig:OGSSeismicity}
\end{figure}

\begin{figure}[h!]
  \centering
  \includegraphics[width=\textwidth]{OGSCatalog/GaMMA0.1/img/OGSAvailability.png}
  \caption{Stations availability during the observation period.}
  \label{fig:OGSAvailability}
\end{figure}

This target dataset provides a challenging testbed due to the region's moderate 
seismicity, crustal heterogeneity, and the presence of significant 
anthropogenic noise sources, making it ideal for assessing model robustness.

To compare the \ac{ML} results with the OGS reference catalog, we relocated the
OGS events using the 1D NonLinLoc (\cite{NonLinLoc}) configuration adopted in
this study. The relocated catalog contains 380 events within the study region,
with local magnitudes ranging from $-0.2$ to 4.6. A total of 7,129 phase 
picks were associated with these events (4,126~P and 3,003~S).

\newpage
\newpage
\section{Benchmark Datasets}

The pre-trained models available in SeisBench employed in this work were 
selected according to three principal criteria. First, each dataset must 
provide a large corpus of three-component seismic waveforms with expert-labeled 
P- and S-wave arrival times, ensuring sufficient supervised signal for model 
training. Second, the datasets should include local earthquakes ($\leq 
150$\,km). Third, the recordings should originate from a diverse set of seismic 
stations and tectonic settings, so that the trained models can generalize 
beyond the conditions of any single network. The following subsections describe 
each dataset in detail.

\subsection{INSTANCE}
\label{sec:INSTANCE}

The INSTANCE dataset (\cite{INSTANCE}) constitutes one of the most 
comprehensive regional seismic collections available for phase-picking 
research, comprising nearly 1.3 million three-component waveforms recorded 
across Italy. Unlike global datasets curated for high signal-to-noise ratios, 
INSTANCE preserves the ``natural'' noise environment encountered in operational 
monitoring networks. This includes urban and industrial noise, cultural noise 
from transportation corridors, meteorological effects, and small-magnitude 
events with emergent onsets. This realism makes INSTANCE a particularly 
valuable benchmark for validation. It captures the operational complexity of 
modern monitoring systems, especially in regions where distinguishing weak 
events from anthropogenic noise is critical. The dataset spans a wide range of 
magnitudes and distances, providing a robust basis for training models that 
must generalize to the specific noise characteristics.

\begin{figure}[h!]
\centering
\includegraphics[width=0.8\textwidth]{INSTANCE.png}
\caption{Geographical distribution of the INSTANCE dataset. Image taken from 
\cite{SeisBench}}
\label{fig:INSTANCE_sample}
\end{figure}

\subsection{STEAD}
\label{sec:STEAD}

The STanford EArthquake Dataset (STEAD) (\cite{STEAD}) provides a complementary 
global perspective with more than 1.2 million labeled waveform segments derived 
from diverse tectonic environments. Designed explicitly for machine-learning
research, STEAD is characterized by:
\begin{itemize}
  \item \textbf{High signal clarity}: Rigorous quality control ensures high 
  signal-to-noise ratios for positive examples.
  \item \textbf{Meticulous labeling}: Precise manual and automatic picks for
  both P and S waves.
  \item \textbf{Noise separation}: A distinct subset of noise-only examples to 
  improve model discrimination.
  \item \textbf{Instrument diversity}: Data recorded by a wide variety of 
  sensors across global networks.
\end{itemize}

Its global coverage minimizes overfitting to local path effects and enables the 
development of models capable of transferring across regions, a crucial 
requirement for universal phase pickers.

\begin{figure}[h!]
\centering
\includegraphics[width=0.8\textwidth]{STEAD.png}
\caption{Global distribution of stations represented in the STEAD dataset. 
Image taken from \cite{SeisBench}}
\label{fig:STEAD_sample}
\end{figure}

\subsection{SCEDC}
\label{sec:SCEDC}

The Southern California Earthquake Data Center (SCEDC) dataset (\cite{SCEDC}) 
offers an unmatched volume of seismic recordings, with more than 8.1 million 
waveform samples. Sourced from the Southern California Seismic Network (SCSN), 
this dataset represents decades of high-quality monitoring in one of the most 
seismically active and well-instrumented regions in the world.

The SCEDC dataset is distinguished by its:
\begin{itemize}
    \item \textbf{Data Volume}: Ideally suited for deep learning, the massive 
    size of the dataset allows for the training of high-capacity models with 
    minimal risk of overfitting.
    \item \textbf{Density}: The high density of stations in Southern California 
    provides redundancy and captures detailed wavefield properties.
    \item \textbf{Catalog Completeness}: The associated catalog is highly 
    complete down to very small magnitudes, providing a rigorous ground truth 
    for evaluating detection sensitivity.
\end{itemize}

In this study, SCEDC serves as a benchmark for large-scale training efficiency 
and provides a baseline for model performance on dense, local networks.

\begin{figure}[h!]
\centering
\includegraphics[width=0.8\textwidth]{SCEDC.png}
\caption{Station map of the SCEDC dataset. Image taken from \cite{SeisBench}}
\label{fig:SCEDC_sample}
\end{figure}

\subsection{Original PhaseNet Corpus}
\label{sec:Original}
The original PhaseNet corpus (\cite{PhaseNet}) consists of approximately 
780,000 waveform segments recorded by the Northern California Seismic Network 
(NCSN). This dataset was curated to train the original PhaseNet model and is 
characterized by:
\begin{itemize}
    \item \textbf{Regional Focus}: Data from a single, well-instrumented 
    region, which may limit generalization but provides a consistent training
    environment.
    \item \textbf{Moderate Size}: Smaller than the other datasets, which may
    lead to faster training but potentially less robust models.
    \item \textbf{High-Quality Picks}: The dataset includes high-quality manual
    picks, which are essential for supervised learning.
\end{itemize}
While not as large or diverse as the other datasets, the original PhaseNet
corpus provides a valuable reference point for evaluating the impact of training
data diversity and volume on model performance.

\subsection{Original EQTransformer Corpus}
\label{sec:Original_EQT}
The original EQTransformer corpus (\cite{EQTransformer}) is a subset of the 
STEAD dataset, consisting of approximately 1.2 million waveform segments. This
dataset was specifically curated to train the original EQTransformer model and
is characterized by:
\begin{itemize}
  \item \textbf{Global Coverage}: Data from a wide range of tectonic settings
  and station types, which enhances the potential for model generalization.
  \item \textbf{Large Size}: The extensive volume of data allows for training
  high-capacity models, though it may require significant computational resources.
  \item \textbf{High-Quality Picks}: The dataset includes meticulously labeled
  picks for both P and S waves, which are crucial for supervised learning.
\end{itemize}
The original EQTransformer corpus serves as a benchmark for evaluating the
impact of training data diversity and volume on the performance of
transformer-based models for seismic phase picking.

Table \ref{tab:datasets} summarizes the principal datasets used in this study. 
The selected collections span regional to global scales and encompass curated, 
high-quality annotations essential for training robust deep-learning models. All 
datasets share a common sampling rate of 100~Hz, which is also the target 
sampling rate adopted for preprocessing in this work.

\begin{table}[h!]
  \centering
  \caption{Summary of training data sets used for PhaseNet \cite{PhaseNet} and 
  EQTransformer \cite{EQTransformer} models. Traces: total number of traces; 
  P: number of P-wave samples; S: number of S-wave samples; Events: number of 
  seismic events; Noise: number of noise samples; Region: geographical
  coverage; Dist: distance range (L: local $\leq$ 150 km, R: regional 150 km 
  $\leq\Delta\leq$ 600 km) (\cite{SeisBench}).
  }
  \begin{tabular}{c|c|c|c|c|c|c|c|}
    Data Set &   Traces  & P         &   S       &  Events &   Noise & Region & Dist \\\hline
    INSTANCE & 1,291,537 & 1,159,249 &   713,883 &  54,008 & 132,288 & Italy  &  L/R \\\hline
    STEAD    & 1,265,657 & 1,030,231 & 1,030,231 & 441,705 & 235,426 & Global &  L/R \\\hline
    SCEDC    & 8,111,060 & 7,571,970 & 4,364,155 & 378,528 &       0 & US     &   L  \\\hline
    Original &   779,514 & \multicolumn{2}{c|}{$\ast$} &  234,117 &       0 & US     &   L  \\\hline
    \multicolumn{8}{l}{\footnotesize $\ast$ P and S sample counts not separately reported in (\cite{PhaseNet}); both phase types are present.}
  \end{tabular}
  \label{tab:datasets}
\end{table}